\section{Limitations and next steps}
\label{sec:limitations}

\paragraph{Price panel survivorship, measured.} The price vendor
serves current symbols only, so the panel is conditioned on survival
to the download date: delisted names are absent for their
\emph{entire} history (nothing ``dies'' inside a backtest window ---
we verified that the return construction never drops a live name
mid-window). We measured the incidence and its signal alignment
rather than asserting a direction: among universe-eligible
instruments ($\ge$5 holders), unmapped instruments --- intentionally
excluded ETFs, share classes, delisted and unresolved names --- are
$\sim$80\% of line items but $\sim$33\% of value, and their NCB vote
rate differs from mapped names by $+0.002$ ($0.0223$ vs $0.0201$,
$t=+2.38$): a $\sim$10\% relative tilt in vote intensity, small but
not zero, and the honest bound on how signal-selected the tradable
universe is. The economic direction of the pure-delisting component
cuts \emph{against} the measured spread: bankruptcies concentrate in
the no-vote quintile (their absence flatters Q1) while acquisition
targets concentrate among conviction names and exit at premiums
(their absence penalizes Q5). CRSP-tier data with
delisting-inclusive returns resolves both truncations and lifts
value coverage from 85.4\% to $\sim$98\% --- the single
highest-value data upgrade.

\paragraph{Coverage residuals.} The identifier-service tail
($\sim$2--3k small CUSIPs) and a handful of throttled tickers are
pending ($\to\sim$93\% free); delisted names are structural until
CRSP.

\paragraph{Sample.} Fifty quarters binds everything inferential: the
IPCA alpha test is underpowered at this $T$; the funding-beta design
has two or three informative episodes; the FSS regime finding rests
on a twenty-quarter window; the hyperparameter surfaces are one draw
of history each.

\paragraph{Corporate actions.} Splits are handled by the modal-atom
detector; other corporate actions (mergers, spin-offs, symbol
recycling) are handled conservatively by fail-closed exclusion
rather than event-level treatment --- a measured 0.22\%
contamination of exit flow, documented rather than repaired.

\paragraph{Book visibility.} 13F books are long-only quarterly
photographs: no shorts, cash, options (excluded by type), or
intra-quarter turnover. The frozen-book return understates
dispersion; AUM-based groupings carry leakage.

\paragraph{Next steps, ranked.} (1) CRSP migration (coverage ceiling
and delisting-inclusive shorts, with corporate-event unification via
permno); (2) higher-frequency flow and performance data ---
monthly N-PORT portfolios and Morningstar-tier reported TNA,
subscriptions/redemptions, and returns --- feeding a probabilistic
position \emph{nowcast} between filings (gated on beating the
frozen-book baseline at holdings reconstruction) and re-opening flow
\emph{prediction} with the instrument the literature actually uses
(reported flows; implied 13F flow has no forecastable structure),
with NMF-cluster hierarchical shrinkage for per-manager
sensitivities; (3) accounting fundamentals and a risk model for
characteristic-complete residualization and joint false-discovery
control; (4) short-interest data to sharpen the SSI instrument and
squeeze guard; (5) a second funding cycle to power the funding-beta
design.
